\label{sec:ObliviousRAM}
The following definitions and notations are as defined by Pinkas et al. \cite{pinkas2010oblivious}.
An \textit{Oblivious Random Access Memory (ORAM)} allows the client to hide her access patterns on the storage system.
Hiding the access patterns makes it \textit{oblivious}, which means that any equal-length sequence of clients' data requests to the server are equivalent from the point of view of the eavesdropper.
This eavesdropper might be the server itself.
Thus, the server must only know the number of queries in the sequence.

Therefore, security in context of ORAMs is defined as follows \cite{pinkas2010oblivious}:
\begin{quotation}
    An access pattern $A(y)$ is a sequence of accesses to the server.
    It contains both the indices accessed in the system and the data items read or written.
    An ORAM system is considered \textit{secure} if for any two inputs $y,y'$ of equal length of the client, the access patterns $A(y),A(y')$ are computationally indistinguishable for anyone but the client.
\end{quotation}
